\documentclass[11pt]{article}
% ===== ?⑦궎吏€ =====
\usepackage[utf8]{inputenc}
\usepackage[T1]{fontenc}
\usepackage{amsmath, amssymb, amsthm}
\usepackage{mathrsfs}
\usepackage[margin=1in]{geometry}
\usepackage{hyperref}
\usepackage{cite}
\usepackage{lettrine}
% ===== ?뺣━ ?섍꼍 =====
\theoremstyle{definition}
\newtheorem{theorem}{Theorem}[section]
\newtheorem{lemma}[theorem]{Lemma}
\newtheorem{proposition}[theorem]{Proposition}
\newtheorem{corollary}[theorem]{Corollary}
\newtheorem{conjecture}[theorem]{Conjecture}

\newtheorem{definition}[theorem]{Definition}
\newtheorem{example}[theorem]{Example}
\newtheorem{remark}[theorem]{Remark}

\numberwithin{equation}{section}
\newcommand{\myabstract}{
\noindent
\textsc{Abstract.}
}
% ===== ?섑븰 湲고샇 ?⑥텞??=====
\newcommand{\N}{\mathbb{N}}
\newcommand{\Z}{\mathbb{Z}}
\newcommand{\R}{\mathbb{R}}
\newcommand{\C}{\mathbb{C}}
\newcommand{\eps}{\varepsilon}
\newcommand{\Real}{\operatorname{Re}}

% ===== ?쒕ぉ / ?€??=====
\title{\textbf{A New Decomposition of the Divisor Summatory Function and Its Asymptotic Formula}}
\author{GUEHERN HWANG \\ Department of Mathematics, Ajou University}
\date{\today}

\begin{document}

\maketitle

% ===== Abstract =====
\begin{quote}
\noindent
\textsc{Abstract.}
In this paper, we study the asymptotic behavior of the divisor summatory
function $D(x)$. We establish the exact identity
\begin{equation*}
    D(x) = 1 + \sum_{i \geq 1} (i+1) \pi(x^{1/i})
         + \sum_{j \geq 2} (-1)^j (j-1) T_j(x),
\end{equation*}
which decomposes $D(x)$ via inclusion--exclusion according to the number
of distinct prime factors of $n$. The three terms correspond respectively
to $n = 1$, prime powers, and integers with $\omega(n) \geq 2$, where the
last sum constitutes the main term.

As the first step toward analyzing $D(x)$, we establish an asymptotic
formula with an explicit error term for $T_j(x)$ with fixed $j \geq 2$.
We consider the two-variable Dirichlet series $F(s, y)$ and obtain the
factorization $F(s, y) = \zeta(s)^{2+2y} H(s, y)$, where $H$ is
holomorphic near $s = 1$. Applying the Selberg--Delange
method---a truncated Perron formula, contour deformation, and local
Hankel analysis---we prove that
\begin{equation*}
    T_j(x) \sim \frac{2^j}{j!} \, x \log x \, (\log \log x)^j
    \quad \text{as } x \to \infty.
\end{equation*}
We hope to extend our results to obtain an asymptotic formula for
$D(x)$ itself.
\end{quote}

% ===== Introduction =====
\section{Introduction}

\subsection{The Dirichlet divisor problem}

Let $d(n)$ denote the number of divisors of a positive integer $n$, and let
\begin{equation}
    D(x) := \sum_{n \leq x} d(n)
\end{equation}
be its summatory function. Dirichlet\cite{} proved in 1849 that
\begin{equation}
    D(x) = x \log x + (2\gamma - 1) x + \Delta(x),
\end{equation}
where $\gamma$ is the Euler--Mascheroni constant and
$\Delta(x) = O(x^{1/2})$. The Dirichlet divisor problem asks for the
infimum $\theta$ of all real numbers such that
$\Delta(x) = O(x^{\theta + \eps})$ for every $\eps > 0$.

Hardy \cite{} proved in 1916 that $\theta \geq 1/4$, and it is
conjectured that $\theta = 1/4$. In the other direction,
Vorono\"i \cite{} first improved Dirichlet's bound to $\theta \leq 1/3$
in 1904, and subsequent refinements based on Vinogradov's method of
exponential sums \cite{} led to the current best estimate, due to
Huxley \cite{} in 2003:
\begin{equation}
    \theta \leq 131/416.
\end{equation}

In this paper, rather than pursuing further improvements of $\theta$,
we explore an alternative approach: an exact combinatorial decomposition
of $D(x)$ via inclusion--exclusion on the number of distinct prime
factors of $n$.

\subsection{A new decomposition}

For a positive integer $n$, let $\omega(n)$ denote the number of distinct prime factors of $n$. For each integer $j\geq2$, define
\begin{equation}\label{eq:Tj-def}
    T_j(x) := \sum_{n \leq x}d(n)\binom{\omega(n)}{j}
\end{equation}
where $\binom{\omega(n)}{j}$ is the binomial coefficient. The main identity of this paper is the following.

\begin{theorem}\label{thm:identity}
    For every real number $x\geq1$,
    \begin{equation}\label{eq:main-identity}
        D(x)=1+\sum_{i\geq 1}(i+1)\pi(x^{1/i})+\sum_{j\geq 2}(-1)^j(j-1)T_j(x)
    \end{equation}
    where $\pi(x)$ denotes the prime-counting function.
\end{theorem}

The three terms on the right-hand side of \eqref{eq:main-identity} correspond, respectively, to $n = 1$, prime powers $\omega(n) = 1$, and integers with $\omega(n)\geq 2$. The proof of Theorem~\ref{thm:identity} is given in Section~\ref{sec:identity}.

Classical results of Dirichlet give the asymptotic
$D(x) \sim x \log x$. By the prime number theorem, the prime-power
contribution satisfies
\begin{equation}
     \sum_{i \geq 1} (i+1) \pi(x^{1/i}) \sim \frac{2x}{\log x},
\end{equation}
which is of strictly smaller order than $x \log x$. Hence the
first two terms in the identity are negligible compared with the
main asymptotic of $D(x)$, and the third sum
\begin{equation}
    \sum_{j \geq 2} (-1)^j (j-1) T_j(x)
\end{equation}
must constitute the main term of $D(x)$. Individually,
each $T_j(x)$ grows rapidly with $j$, but the alternating coefficients
$(-1)^j (j-1)$ induce substantial cancellation, reducing the total
contribution to the order $x \log x$. Understanding the
asymptotic behavior of each $T_j(x)$ is therefore the first step toward
an asymptotic formula for $D(x)$ itself.

\begin{example}
    At $x=10000$, direct computation gives
    \begin{align*}
        T_2(10000) &= 297{,}078 & T_3(10000) &= 140{,}640 \\
        T_4(10000) &= 26{,}674  & T_5(10000) &= 1{,}216
    \end{align*}
    and $T_j(10000)=0$ for $j\geq6$. The prime-power contribution equals
$2{,}711$, and the identity \eqref{eq:main-identity} yields
    \begin{align*}
        1+2{,}711+(T_2-2T_3+3T_4-4T_5)=1+2{,}711+90{,}956=93{,}668=D(10000),
    \end{align*}
    confirming Theorem~\ref{thm:identity} numerically. Moreover, the ratio between
$T_j(x)$ and its conjectured asymptotic $\frac{2^j}{j!}x\log x(\log\log x)^j$ converges to 1 only very slowly: at $x=10^9$, this ratio is 0.43 for
$j=2$ and 0.21 for $j=3$. This slow convergence motivates the refined
asymptotic expansion of $T_j(x)$ developed in Section~\ref{sec:extraction}.
\end{example}

\subsection{Main results}

To analyze $T_j(x)$, we consider the two-variable Dirichlet series
\begin{equation}
    F(s, y) := \sum_{n \geq 1} \frac{d(n)(1+y)^{\omega(n)}}{n^s},
\end{equation}
which encodes $T_j(x)$ as the coefficient of $y^j$ in the generating
relation $\sum_{j \geq 0} T_j(x) y^j$. Here, we extend the definition
of $T_j(x)$ to all $j \geq 0$ by the formula~\eqref{eq:Tj-def}, so
that $T_0(x) = D(x)$ and $T_1(x) = \sum_{n \leq x} d(n) \omega(n)$.
The series $F(s, y)$ admits the factorization
\begin{equation}
    F(s, y) = \zeta(s)^{2+2y} H(s, y),
\end{equation}
where $H(s, y)$ is holomorphic in $s$ on a neighborhood of
$\Real(s) \geq 1$ and uniformly bounded for $y$ in a small disk
around $0$. This factorization isolates the singularity of
$F(s, y)$ at $s = 1$, and it is the starting point of the
Selberg--Delange method.

Our main asymptotic result is the following.

\begin{theorem}\label{thm:main}
For each fixed integer $j \geq 2$, as $x \to \infty$,
\begin{equation}
\begin{aligned}
    T_j(x)
    &= \frac{2^j}{j!} \, x \log x \, (\log \log x)^j \\
    &\quad + \frac{2^{j-1}}{(j-1)!}
       \bigl(\partial_y H(1, 0) - 2(1 - \gamma)\bigr)
       x \log x \, (\log \log x)^{j-1} \\
    &\quad + O\bigl(x \log x \, (\log \log x)^{j-2}\bigr).
\end{aligned}
\end{equation}
In particular,
\begin{equation}
    T_j(x) \sim \frac{2^j}{j!} \, x \log x \, (\log \log x)^j.
\end{equation}
\end{theorem}

The proof of Theorem~\ref{thm:main} is given in Sections~\ref{sec:perron}--\ref{sec:extraction}. The essential ingredients
are the truncated Perron formula, contour deformation around
$s = 1$, and a local Hankel analysis in the $w$-plane obtained
via the substitution $w = (s - 1) \log x$. The coefficient $\partial_y H(1, 0)$ admits an explicit expression as an absolutely convergent series over primes; we record this
expression in Section~\ref{sec:extraction}, where it is used in the proof of Theorem~\ref{thm:main}.

Combining Theorem~\ref{thm:main} with the identity~\eqref{eq:main-identity},
we hope to obtain an asymptotic formula for $D(x)$ itself. We
discuss this direction in Section~\ref{sec:outlook}.

\subsection{Outline of the paper}

The paper is organized as follows. In Section~\ref{sec:identity},
we prove the identity (Theorem~\ref{thm:identity}) by partitioning
$D(x)$ according to $\omega(n)$ and applying an inclusion--exclusion
argument to the composite part. In Section~\ref{sec:dirichlet-series},
we introduce the two-variable Dirichlet series $F(s, y)$ and establish
its factorization $F(s, y) = \zeta(s)^{2+2y} H(s, y)$, together with
the analytic properties of $H(s, y)$ needed in the sequel.

Sections~\ref{sec:perron}--\ref{sec:extraction} are devoted to the
proof of Theorem~\ref{thm:main}. In Section~\ref{sec:perron}, we
apply the truncated Perron formula to the generating partial sum
$A(x, y) := \sum_{n \leq x} d(n) (1+y)^{\omega(n)}$ and deform the
contour around $s = 1$. In Section~\ref{sec:hankel}, we carry out
the local Hankel analysis via the substitution $w = (s-1) \log x$
and evaluate the main and lower-order contributions. In
Section~\ref{sec:extraction}, we extract the asymptotic formula for
$T_j(x)$ from that of $A(x, y)$ by differentiating in $y$, thereby
completing the proof of Theorem~\ref{thm:main}.

Finally, in Section~\ref{sec:outlook}, we combine
Theorem~\ref{thm:main} with the identity~\eqref{eq:main-identity}
and discuss the remaining obstacles toward an asymptotic formula
for $D(x)$.

% ===== Section 2 =====
\section{Proof of the identity}\label{sec:identity}

In this section, we prove the identity stated in
Theorem~\ref{thm:identity}. The proof proceeds by partitioning
the sum $D(x) = \sum_{n \leq x} d(n)$ according to the value of
$\omega(n)$, and then applying an inclusion--exclusion argument
to the composite part.

\begin{proof}[Proof of Theorem~\ref{thm:identity}]
Partition the integers $n$ with $1 \leq n \leq x$ into three
disjoint classes:
\begin{itemize}
    \item[(i)] $n = 1$ (equivalently, $\omega(n) = 0$);
    \item[(ii)] prime powers, i.e., $\omega(n) = 1$;
    \item[(iii)] composite integers with $\omega(n) \geq 2$.
\end{itemize}
Writing $D(x)$ as the sum of contributions from each class, we
have
\begin{equation}
    D(x) = d(1) + \sum_{\substack{n \leq x \\ \omega(n) = 1}} d(n)
               + \sum_{\substack{n \leq x \\ \omega(n) \geq 2}} d(n).
\end{equation}

Since $d(1) = 1$, the first term equals $1$.

For the second term, every integer $n$ with $\omega(n) = 1$ is of
the form $n = p^i$ for a unique prime $p$ and a unique integer
$i \geq 1$, and $d(p^i) = i + 1$. The number of prime powers
$p^i \leq x$ with exponent exactly $i$ equals $\pi(x^{1/i})$.
Hence
\begin{equation}
    \sum_{\substack{n \leq x \\ \omega(n) = 1}} d(n)
    = \sum_{i \geq 1} (i+1) \pi(x^{1/i}).
\end{equation}

For the third term, we apply inclusion--exclusion on the set of
distinct prime divisors of $n$. For each integer $j \geq 2$,
recall the definition
\[
    T_j(x) = \sum_{n \leq x} d(n) \binom{\omega(n)}{j}.
\]
We claim that
\begin{equation}
    \sum_{\substack{n \leq x \\ \omega(n) \geq 2}} d(n)
    = \sum_{j \geq 2} (-1)^j (j-1) T_j(x).
\end{equation}

To prove~(2.3), interchange the order of summation on the
right-hand side:
\[
    \sum_{j \geq 2} (-1)^j (j-1) T_j(x)
    = \sum_{n \leq x} d(n) \sum_{j \geq 2} (-1)^j (j-1)
      \binom{\omega(n)}{j}.
\]
For $n$ with $\omega(n) \in \{0, 1\}$, we have
$\binom{\omega(n)}{j} = 0$ for all $j \geq 2$, so such $n$ do not
contribute. For $n$ with $\omega(n) = m \geq 2$, we use the
following combinatorial identity.

\begin{lemma}
For every integer $m \geq 2$,
\begin{equation}
    \sum_{j=2}^{m} (-1)^j (j-1) \binom{m}{j} = 1.
\end{equation}
\end{lemma}

\begin{proof}[Proof of the lemma]
Starting from the standard identities
\[
    \sum_{j=0}^{m} (-1)^j \binom{m}{j} = 0
    \quad \text{and} \quad
    \sum_{j=0}^{m} (-1)^j \, j \binom{m}{j} =\left.\frac{d}{dx}(1-x)^m\right|_{x=1}=0 \ \text{(for } m \geq 2),
\]
we combine them to obtain
\[
    \sum_{j=0}^{m} (-1)^j (j-1) \binom{m}{j} = 0
    \quad \text{(for } m \geq 2).
\]
Separating the $j = 0$ and $j = 1$ terms:
\[
    1 \cdot (-1) + 0 \cdot \binom{m}{1}
    + \sum_{j=2}^{m} (-1)^j (j-1) \binom{m}{j} = 0,
\]
which yields
\[
    \sum_{j=2}^{m} (-1)^j (j-1) \binom{m}{j} = 1. \qedhere
\]
\end{proof}

Applying the lemma, we obtain
\[
    \sum_{n \leq x} d(n) \sum_{j \geq 2} (-1)^j (j-1)
    \binom{\omega(n)}{j}
    = \sum_{\substack{n \leq x \\ \omega(n) \geq 2}} d(n) \cdot 1,
\]
which establishes~(2.3).

Combining the contributions from all three classes yields
\[
    D(x) = 1 + \sum_{i \geq 1} (i+1) \pi(x^{1/i})
             + \sum_{j \geq 2} (-1)^j (j-1) T_j(x),
\]
completing the proof.
\end{proof}

\section{Dirichlet series and factorization}\label{sec:dirichlet-series}

In this section, we introduce the two-variable Dirichlet series
$F(s, y)$ and establish its key analytic properties. The main
output is the factorization
$F(s, y) = \zeta(s)^{2+2y} H(s, y)$, where $H(s, y)$ is holomorphic
in a neighborhood of $s = 1$ and uniformly bounded in a small disk
around $y = 0$.

\subsection{Euler product}\label{sec:euler-product}

Recall the definition
\begin{equation}
    F(s, y) := \sum_{n \geq 1} \frac{d(n)(1+y)^{\omega(n)}}{n^s}.
\end{equation}
Since the arithmetic function $n \mapsto d(n)(1+y)^{\omega(n)}$ is
multiplicative in $n$ for each fixed $y$, the series admits an Euler
product representation for $\Real(s) > 1$.

\begin{lemma}\label{lem:euler-product}
For $\Real(s) > 1$ and every $y \in \C$,
\begin{equation}\label{eq:euler-product}
    F(s, y) = \prod_p \left(1 + (1+y)\left((1 - p^{-s})^{-2} - 1\right)\right),
\end{equation}
where the product is over all primes $p$.
\end{lemma}

\begin{proof}
By multiplicativity, for $\Real(s) > 1$ we have
\[
    F(s, y) =\sum_{n \geq 1}\frac{d(n)(1+y)^{\omega(n)}}{n^s} = \prod_p \left(\sum_{k \geq 0}
    \frac{d(p^k)(1+y)^{\omega(p^k)}}{p^{ks}}\right).
\]
For $k = 0$, the summand equals $1$. For $k \geq 1$, we have
$d(p^k) = k + 1$ and $\omega(p^k) = 1$, so the local factor becomes
\[
    1 + (1+y) \sum_{k \geq 1} \frac{k+1}{p^{ks}}.
\]
Setting $u := p^{-s}$ and using the standard identity
\[
    \sum_{k \geq 1} (k+1) u^k = \frac{1}{(1-u)^2} - 1
    \quad \text{for } |u| < 1,
\]
the local factor equals $1 + (1+y)\bigl((1-p^{-s})^{-2} - 1\bigr)$,
which proves~\eqref{eq:euler-product}.
\end{proof}

\subsection{Factorization isolating the singularity}\label{sec:factorization}

The Riemann zeta function admits the Euler product
$\zeta(s) = \prod_p (1 - p^{-s})^{-1}$ for $\Real(s) > 1$, so that
$\zeta(s)^{2+2y} = \prod_p (1 - p^{-s})^{-(2+2y)}$ on the same
half-plane, provided a branch of the logarithm is fixed. We use the
principal branch, which is analytic in $s$ in a neighborhood of
any point with $\Real(s) > 1$ and in $y$ in a neighborhood of $0$.

Define
\begin{equation}\label{eq:H-def}
    H(s, y) := \prod_p (1 - p^{-s})^{2+2y}
    \left(1 + (1+y)\left((1 - p^{-s})^{-2} - 1\right)\right).
\end{equation}

\begin{theorem}\label{thm:factorization}
For $\Real(s) > 1$ and $y$ in a neighborhood of $0$,
\begin{equation}\label{eq:factorization}
    F(s, y) = \zeta(s)^{2+2y} H(s, y).
\end{equation}
\end{theorem}

\begin{proof}
This follows by combining Lemma~\ref{lem:euler-product} with the
Euler product for $\zeta(s)^{2+2y}$ and regrouping factors.
\end{proof}

The purpose of the factorization~\eqref{eq:factorization} is to
isolate the singular behavior of $F(s, y)$ at $s = 1$ into the
factor $\zeta(s)^{2+2y}$, leaving $H(s, y)$ as a well-behaved
function in a larger region, as we now establish.

\subsection{Analytic properties of \texorpdfstring{$H(s, y)$}{H(s,y)}}\label{sec:H-analytic}

We show that the Euler product~\eqref{eq:H-def} defines a
holomorphic function on a strip containing $\Real(s) = 1$,
uniformly for $y$ in a small disk around $0$. The key is that
each local factor $H_p(s, y)$ deviates from $1$ by a quantity of
order $p^{-2\Real(s)}$.

\begin{lemma}\label{lem:local-factor-estimate}
Fix $\eta \in (0, 1/2)$ and $\delta \in (0, 1/2)$. Write
$u := p^{-s}$ and denote the local factor by
\[
    H_p(s, y) := (1 - u)^{2+2y}\left(1 + (1+y)\left((1-u)^{-2} - 1\right)\right).
\]
Then for $|y| \leq \eta$ and $\Real(s) \geq 1 - \delta$,
\begin{equation}\label{eq:local-factor-estimate}
    H_p(s, y) = 1 + O_{\eta, \delta}(u^2),
\end{equation}
where the implied constant depends only on $\eta$ and $\delta$.
\end{lemma}

\begin{proof}
The condition $\Real(s) \geq 1 - \delta$ implies that, for every
prime $p \geq 2$,
\[
    |u| = p^{-\Real(s)} \leq p^{-(1 - \delta)} \leq 2^{-(1-\delta)} < 1,
\]
since $\delta < 1/2$. In particular, $|u|$ is bounded away from
$1$, and the power series expansions
\[
    (1 - u)^{-2} - 1 = 2u + 3u^2 + 4u^3 + \cdots = 2u + O_\delta(u^2),
\]
\[
    (1 - u)^{2+2y} = 1 - (2+2y)u + \binom{2+2y}{2}u^2 + \cdots
                   = 1 - (2+2y)u + O_{\eta, \delta}(u^2),
\]
converge absolutely, with remainders uniform for $|y| \leq \eta$
and $\Real(s) \geq 1 - \delta$. Multiplying the expansions:
\[
    H_p(s, y) = \bigl(1 - (2+2y)u + O_{\eta, \delta}(u^2)\bigr)
                \bigl(1 + (1+y)(2u + O_\delta(u^2))\bigr).
\]
The coefficient of $u^1$ on the right is $-(2+2y) + 2(1+y) = 0$.
Hence $H_p(s, y) = 1 + O_{\eta, \delta}(u^2)$,
proving~\eqref{eq:local-factor-estimate}.
\end{proof}

% =====================================================
% Proposition 3.4 (revised): H is jointly holomorphic
% =====================================================

\begin{proposition}\label{prop:H-analytic}
Fix $\eta \in (0, 1/2)$ and $\delta \in (0, 1/2)$, and let
\begin{equation}\label{eq:H-region}
    \mathcal{R}_{\delta, \eta} := \{(s, y) \in \C^2 :
    \Real(s) \geq 1 - \delta, \ |y| \leq \eta\}.
\end{equation}
Then the partial products of the Euler product~\eqref{eq:H-def}
converge absolutely and uniformly on $\mathcal{R}_{\delta, \eta}$
to a function $H(s, y)$. For every $y$ with $|y| \leq \eta$, the
function $H(s, y)$ is holomorphic in $s$ on the open strip
$\Real(s) > 1 - \delta$. Moreover, $H(s, y)$ is jointly holomorphic
in $(s, y)$ on $\operatorname{int}(\mathcal{R}_{\delta, \eta})$,
and $H$ is uniformly bounded on $\mathcal{R}_{\delta, \eta}$.
\end{proposition}

\begin{proof}
By Lemma~\ref{lem:local-factor-estimate}, for every
$(s, y) \in \mathcal{R}_{\delta, \eta}$ and every prime $p$,
\begin{equation}\label{eq:Hp-bound}
    |H_p(s, y) - 1| \leq C_{\eta, \delta} \, p^{-2(1-\delta)},
\end{equation}
where the constant $C_{\eta, \delta}$ depends only on $\eta$ and
$\delta$. Since $\delta < 1/2$, we have $2(1 - \delta) > 1$, and
the series $\sum_n n^{-2(1-\delta)}$ converges; in particular,
$\sum_p p^{-2(1-\delta)}$ converges. By the Weierstrass $M$-test,
the series $\sum_p (H_p(s, y) - 1)$ converges absolutely and
uniformly on $\mathcal{R}_{\delta, \eta}$. By the standard theorem
on infinite products---which states that an infinite product
$\prod_n (1 + u_n)$ converges absolutely if and only if
$\sum_n u_n$ does, and that the convergence of the partial products
is uniform on a set whenever the corresponding series converges
uniformly there---the partial products
$\prod_{p \leq P} H_p(s, y)$ converge absolutely and uniformly on
$\mathcal{R}_{\delta, \eta}$ to a limit, which we denote $H(s, y)$.

To deduce holomorphy in $s$, fix $y_0 \in \C$ with $|y_0| \leq \eta$
and let $s_0 \in \C$ be any point with $\Real(s_0) > 1 - \delta$.
Choose $r > 0$ small enough that $r < \Real(s_0) - (1 - \delta)$,
and define
\begin{equation}\label{eq:K-disk}
    K := \overline{D}(s_0, r)
       = \{s \in \C : |s - s_0| \leq r\}.
\end{equation}
The set $K$ is closed and bounded in $\C$, so by the Heine--Borel
theorem it is compact. By the choice of $r$, every $s \in K$
satisfies
\[
    \Real(s) \geq \Real(s_0) - r > 1 - \delta,
\]
and hence $K \times \{y_0\} \subset \mathcal{R}_{\delta, \eta}$.
The uniform convergence of the partial products on
$\mathcal{R}_{\delta, \eta}$ established above therefore restricts
to uniform convergence of $\prod_{p \leq P} H_p(s, y_0)$ on $K$
as $P \to \infty$.

Each finite partial product $\prod_{p \leq P} H_p(s, y_0)$ is
holomorphic in $s$ on the open disk
$D(s_0, r) = \operatorname{int}(K)$, since each factor
$H_p(\cdot, y_0)$ is. By the theorem on uniform limits of
holomorphic functions (Weierstrass), the limit $H(\cdot, y_0)$ is
holomorphic on $D(s_0, r)$, and in particular at $s_0$. Since
$s_0$ with $\Real(s_0) > 1 - \delta$ was arbitrary, $H(\cdot, y_0)$
is holomorphic on the open strip $\Real(s) > 1 - \delta$. Since
$y_0$ with $|y_0| \leq \eta$ was also arbitrary, the holomorphy
claim in $s$ is proved.

For joint holomorphy, each finite partial product
$\prod_{p \leq P} H_p(s, y)$ is holomorphic in $(s, y)$ on
$\operatorname{int}(\mathcal{R}_{\delta, \eta})$, since each factor
$H_p(s, y)$ is. The uniform convergence of the partial products on
$\mathcal{R}_{\delta, \eta}$ established above restricts to uniform
convergence on every compact subset of
$\operatorname{int}(\mathcal{R}_{\delta, \eta})$, since
$\operatorname{int}(\mathcal{R}_{\delta, \eta}) \subset
\mathcal{R}_{\delta, \eta}$. By the theorem on uniform limits of
holomorphic functions in two complex variables, the limit
$H(s, y)$ is jointly holomorphic on
$\operatorname{int}(\mathcal{R}_{\delta, \eta})$.

Uniform boundedness on $\mathcal{R}_{\delta, \eta}$ follows
from~\eqref{eq:Hp-bound}:
\[
    |H(s, y)|
    \leq \prod_p \bigl(1 + |H_p(s, y) - 1|\bigr)
    \leq \exp\!\left(\sum_p |H_p(s, y) - 1|\right)
    \leq \exp\!\left(C_{\eta, \delta} \sum_p p^{-2(1-\delta)}\right),
\]
the right-hand side being finite and independent of
$(s, y) \in \mathcal{R}_{\delta, \eta}$.
\end{proof}


% =====================================================
% Proposition 3.5 (new): H(s, 0) = 1
% =====================================================

\begin{proposition}\label{prop:H-at-zero}
The function $H(s, y)$ satisfies
\begin{equation}\label{eq:H-at-y-zero}
    H(s, 0) \equiv 1 \quad \text{on } \Real(s) > 1/2.
\end{equation}
In particular, $H(1, 0) = 1$.
\end{proposition}

\begin{proof}
On $\Real(s) > 1$, the factorization~\eqref{eq:factorization}
specialized at $y = 0$ reads
\[
    F(s, 0) = \zeta(s)^2 \, H(s, 0).
\]
On the other hand, the Dirichlet series defining $F$ gives
\[
    F(s, 0) = \sum_{n \geq 1} \frac{d(n)}{n^s} = \zeta(s)^2
    \quad \text{on } \Real(s) > 1.
\]
Comparing the two expressions yields
$\zeta(s)^2 H(s, 0) = \zeta(s)^2$ on $\Real(s) > 1$. Since
$\zeta(s) \neq 0$ on $\Real(s) > 1$ (which follows directly from
the Euler product representation of $\zeta$), we may divide both
sides by $\zeta(s)^2$ to obtain
\[
    H(s, 0) \equiv 1 \quad \text{on } \Real(s) > 1.
\]

By Proposition~\ref{prop:H-analytic}, applied with any
$\delta \in (0, 1/2)$, the function $H(s, 0)$ is holomorphic on
$\Real(s) > 1 - \delta$; letting $\delta \to 1/2$ shows that
$H(s, 0)$ is holomorphic on the connected open set
$\Real(s) > 1/2$. The constant function $1$ is also holomorphic
there, and the two functions agree on the open subset
$\Real(s) > 1$. By the identity theorem, they agree on all of
$\Real(s) > 1/2$. In particular, $H(1, 0) = 1$.
\end{proof}

\section{Perron formula and contour deformation}\label{sec:perron}

To extract the asymptotic behavior of $T_j(x)$ from the Dirichlet
series $F(s, y)$, we work with the generating partial sum
\begin{equation}\label{eq:Axy-def}
    A(x, y) := \sum_{n \leq x} d(n) (1 + y)^{\omega(n)},
\end{equation}
which satisfies the generating relation
\begin{equation}\label{eq:Axy-generating}
    A(x, y) = \sum_{j \geq 0} T_j(x) y^j.
\end{equation}
Throughout this section, we fix one parameter:
\begin{itemize}
    \item $\eta \in (0, 1/2)$, controlling the size of the disk
    $|y| \leq \eta$.
\end{itemize}

To avoid threading $T$-dependent constants through every estimate
borrowed from Section~\ref{sec:dirichlet-series}, we also fix once
and for all an auxiliary constant
\begin{equation}\label{eq:delta-prime}
    \delta' := \frac{1}{4},
\end{equation}
and apply the results of Section~\ref{sec:dirichlet-series} on the
fixed region $\mathcal{R}_{\delta', \eta}$. Constants such as
$C_{\eta, \delta'}$ inherited from those results are absolute
(they do not depend on $T$).

A second quantity, $\lambda_T$, will govern the depth of the
contour deformation in Section~\ref{sec:contour-decomposition}.
In contrast to the free parameter $\delta \in (0, 1/2)$ used in
Section~\ref{sec:dirichlet-series} and the fixed constant $\delta'$
just defined, $\lambda_T$ is a function of the truncation height
$T \geq 1$, defined as
\begin{equation}\label{eq:lambda-T}
    \lambda_T := \frac{c_0}{2 \log(T + 2)},
\end{equation}
where $c_0$ is the constant of Theorem~\ref{thm:zero-free}. By construction,
\begin{equation}\label{eq:lambda-T-bound}
    \lambda_T
    \leq \frac{c_0}{2 \log 3}
    < \frac{1}{4 \log 3}
    < \frac{1}{4} = \delta'
    \qquad (T \geq 1),
\end{equation}
and $\lambda_T \to 0$ as $T \to \infty$. Since $\lambda_T < \delta'$,
the contour $\mathcal{C}(T, \rho)$ constructed in
Section~\ref{sec:contour-decomposition} (and the closed loop
$\Gamma$ it forms) lies in $\mathcal{R}_{\delta', \eta}$ for every
$T \geq 1$ and every $|y| \leq \eta$. Consequently, all bounds on
$H$ inherited from Section~\ref{sec:dirichlet-series} are uniform
in $T$ on the contour and its enclosed region.

We also abbreviate $\alpha := 2 + 2y$. Thus the
factorization~\eqref{eq:factorization} becomes
$F(s, y) = \zeta(s)^\alpha H(s, y)$. For $|y| \leq \eta$ with
$\eta < 1/2$, the exponent $\alpha$ ranges in a disk of radius
$2\eta < 1$ around $2$, so in particular $\alpha$ is bounded away
from non-positive integers.

\subsection{The truncated Perron formula}\label{sec:perron-formula}

We quote the following standard form of the truncated Perron
formula (see, e.g., \cite[Theorem II.2.2]{}).

\begin{theorem}[Truncated Perron formula]\label{thm:perron}
Let $f(s) = \sum_{n \geq 1} a_n n^{-s}$ be a Dirichlet series with
abscissa of absolute convergence $\sigma_a$. For real numbers
$c > \sigma_a$, $x \geq 2$, and $T \geq 1$ with $x \notin \mathbb{Z}$,
\begin{equation}\label{eq:perron}
    \sum_{n \leq x} a_n
    = \frac{1}{2\pi i} \int_{c - iT}^{c + iT} f(s) \frac{x^s}{s} \, ds
    + R(x, T),
\end{equation}
where the truncation error satisfies
\begin{equation}\label{eq:perron-error}
    R(x, T) \ll \frac{x^c}{T} \sum_{n \geq 1} \frac{|a_n|}{n^c
    \left(1 + T \left|\log \frac{x}{n}\right|\right)}
    + \frac{x^c}{T \log x}.
\end{equation}
\end{theorem}

We apply Theorem~\ref{thm:perron} to the Dirichlet series
$F(s, y) = \sum_{n \geq 1} d(n) (1+y)^{\omega(n)} n^{-s}$. For
$|y| \leq \eta$ with $\eta$ sufficiently small, the series
converges absolutely for $\Real(s) > 1$, so we may take any $c > 1$.
Choose, for instance, $c = 1 + 1/\log x$, which is the standard
choice to balance the main term size and the truncation error.

\begin{proposition}\label{prop:perron-applied}
For $x \geq 2$ with $x \notin \mathbb{Z}$, $|y| \leq \eta$, and
$T \geq 1$,
\begin{equation}\label{eq:Axy-perron}
    A(x, y) = \frac{1}{2\pi i} \int_{c - iT}^{c + iT} F(s, y)
    \frac{x^s}{s} \, ds + E_{\mathrm{Perron}}(x, T, y),
\end{equation}
where $c = 1 + 1/\log x$ and the truncation error satisfies
\begin{equation}\label{eq:Axy-perron-error}
    E_{\mathrm{Perron}}(x, T, y) \ll_{\eta, \delta'}
    \frac{x (\log x)^{2 + 2\eta}}{T}.
\end{equation}
\end{proposition}

\begin{proof}
Apply Theorem~\ref{thm:perron} with $a_n = d(n)(1+y)^{\omega(n)}$
and $c = 1 + 1/\log x$. Denote the resulting Perron truncation
error by $R(x, T, y)$. Using $|1+y| \leq 1+\eta$ and the trivial
bound $1 + T|\log(x/n)| \geq 1$, the first term
of~\eqref{eq:perron-error} satisfies
\begin{align*}
    \frac{x^c}{T} \sum_{n \geq 1}
    \frac{\left| d(n)(1+y)^{\omega(n)} \right|}
         {n^c \left(1 + T \left|\log \frac{x}{n}\right|\right)}
    \leq \frac{x^c}{T} \sum_{n \geq 1}
    \frac{d(n)(1+ \left|y\right|)^{\omega(n)}}{n^c}
    \leq \frac{x^c}{T} \sum_{n \geq 1}
    \frac{d(n)(1+\eta)^{\omega(n)}}{n^c}
    = \frac{x^c}{T} F(c, \eta).
\end{align*}
Hence
\begin{equation}\label{eq:R-bound-raw}
    R(x, T, y) \ll_{\eta, \delta'} \frac{x^c}{T} F(c, \eta)
    + \frac{x^c}{T \log x}.
\end{equation}

We now estimate $F(c, \eta)$ at $c = 1 + 1/\log x$. By the
factorization~\eqref{eq:factorization},
\[
    F(c, \eta) = \zeta(c)^{2 + 2\eta} H(c, \eta).
\]
Since $\zeta(s)$ has a simple pole at $s = 1$ with residue $1$,
the standard expansion
$\zeta(s) = \frac{1}{s - 1} + \gamma + O(s-1)$ near $s = 1$ gives
\[
    \zeta\!\left(1 + \tfrac{1}{\log x}\right)
    = \log x + \gamma + O\!\left(\tfrac{1}{\log x}\right)
    \sim \log x \quad (x \to \infty).
\]
Hence $\zeta(c)^{2+2\eta} \ll_\eta (\log x)^{2+2\eta}$. Moreover, Proposition~\ref{prop:H-analytic}, applied with
$\delta = \delta'$, implies that $H(c, \eta)$ is bounded by
a constant depending only on $\eta$ and $\delta'$. Therefore
\begin{equation}\label{eq:F-at-c}
    F(c, \eta) \ll_{\eta, \delta'} (\log x)^{2 + 2\eta}.
\end{equation}

Substituting~\eqref{eq:F-at-c} into~\eqref{eq:R-bound-raw} and
using $x^c = x \cdot x^{c - 1} = x \cdot x^{1/\log x} = e\,x$,
\[
    R(x, T, y)
    \ll_{\eta, \delta'} \frac{x (\log x)^{2+2\eta}}{T}
    + \frac{x}{T \log x}
    \ll_{\eta, \delta'} \frac{x (\log x)^{2+2\eta}}{T},
\]
where the second term is absorbed into the first. Setting
$E_{\mathrm{Perron}}(x, T, y) := R(x, T, y)$ proves the proposition.
\end{proof}

The Perron truncation error~\eqref{eq:Axy-perron-error} is of size $x(\log x)^{2+2\eta}/T$. Since $T$ appears in the denominator, this term decreases as $T$ grows, and it can be made negligible compared to the main term by choosing $T$ sufficiently large as a function of $x$. The precise choice of $T = T(x)$ will be made in Section~\ref{sec:contour-estimates}, after combining this error with the contributions from the contour deformation.

\subsection{Analytic preliminaries}\label{sec:analytic-prelim}

To control the contour integrals arising from the deformation in
Section~\ref{sec:contour-decomposition}, we will need:
\begin{itemize}
    \item[(a)] a region in which $\zeta(s)$ does not vanish, so
    that the contour can avoid its zeros;
    \item[(b)] a continuous branch of $\zeta(s)^{2 + 2y}$ on a
    neighborhood of the contour;
    \item[(c)] a uniform bound on $|\zeta(s)|$ along vertical
    lines in the critical strip;
    \item[(d)] a corresponding bound on $|F(s, y)|$ combining
    (b) and (c).
\end{itemize}
We collect these analytic ingredients here, before constructing
the contour itself in Section~\ref{sec:contour-decomposition}.

\subsubsection*{Zero-free region}

For the contour deformation to be valid, we need a region in which
$\zeta(s) \neq 0$. We use the following classical result of
de la Vall\'ee Poussin, quoted without proof; see, e.g.,
\cite[Ch.~3]{} or \cite[\S6.6]{}.

\begin{theorem}[de la Vall\'ee Poussin zero-free region]
\label{thm:zero-free}
There exists an absolute constant $c_0 \in (0, 1/2)$ such that
$\zeta(s) \neq 0$ on the open region
\begin{equation}\label{eq:zero-free-region}
    \mathcal{Z} := \left\{ s = \sigma + it \in \C :
    \sigma > 1 - \frac{c_0}{\log(|t| + 2)} \right\}
    \setminus \{1\}.
\end{equation}
\end{theorem}

We shall use two different strip parameters. The fixed value
$\delta'=1/4$ introduced above is used to keep the Euler product
factor $H(s,y)$ uniformly controlled on a $T$-independent region.
For growth estimates on the deformed contour we use a smaller
parameter $\delta_*$, chosen once and for all so that
\begin{equation}\label{eq:delta-star}
    0 < \delta_* < \min\!\left(\delta', \frac{c_0}{2}\right).
\end{equation}
This separation is useful because the exponent in the vertical
growth bound below tends to $\delta_*$ as $\eta,\varepsilon \to 0^+$,
whereas the contour itself is still allowed to lie in the larger
region $\mathcal{R}_{\delta',\eta}$ where $H$ is controlled.

The boundary of $\mathcal{Z}$ is curved, and integrating along it
would complicate the estimation of the resulting contour pieces.
We will instead use a fixed vertical line lying strictly inside
$\mathcal{Z}$ for all $|t| \leq T$.

\subsubsection*{Continuous branch of \texorpdfstring{$\zeta(s)^\alpha$}{zeta(s)^alpha}}

To make sense of $\zeta(s)^\alpha = \zeta(s)^{2 + 2y}$ on a
neighborhood of the contour, we specify a branch on the region
\begin{equation}\label{eq:D-region}
    \mathcal{D} := \mathcal{Z} \setminus \mathcal{S},
\end{equation}
where $\mathcal{S} := \{s \in \C : \Im(s) = 0,\ \Real(s) \leq 1\}$
is the slit along the real axis to the left of $s = 1$. The slit
$\mathcal{S}$ disconnects the segment $\mathcal{Z} \cap \mathcal{S}$
from the rest of $\mathcal{Z}$ from above and below, rendering
$\mathcal{D}$ simply connected.\footnote{%
Indeed, $\mathcal{Z}$ is a slit-once-punctured domain that is not
simply connected, since a small loop around $s = 1$ inside
$\mathcal{Z}$ is not nullhomotopic. Removing the
entire ray $(-\infty, 1]$ (equivalently, the connected component
$\mathcal{Z} \cap \mathcal{S}$ of $\mathcal{S}$ inside
$\mathcal{Z}$, which contains $s = 1$) cuts every such loop,
making the resulting domain simply connected.}

On $\mathcal{D}$, both $\zeta(s)$ and $\zeta'(s)$ are holomorphic,
and $\zeta(s) \neq 0$. Fix any real number $c_* > 1$ (for
example $c_* = 2$); since $\zeta(c_*) > 0$, the principal
logarithm $\log \zeta(c_*) \in \R$ is well-defined. Define
\begin{equation}\label{eq:L-def}
    L(s) := \log \zeta(c_*)
          + \int_{c_*}^{s} \frac{\zeta'(\omega)}{\zeta(\omega)}\, d\omega
    \qquad (s \in \mathcal{D}),
\end{equation}
where the integral is taken along any path in $\mathcal{D}$ from
$c_*$ to $s$. Since $\mathcal{D}$ is simply connected and
$\zeta'/\zeta$ is holomorphic on $\mathcal{D}$, the integral is
path-independent, and $L$ is a well-defined holomorphic function
on $\mathcal{D}$.

We claim $e^{L(s)} = \zeta(s)$ on $\mathcal{D}$. Indeed,
$L'(s) = \zeta'(s)/\zeta(s)$ on $\mathcal{D}$, so
\[
    \frac{d}{ds}\!\bigl(e^{-L(s)} \zeta(s)\bigr)
    = e^{-L(s)} \bigl(\zeta'(s) - L'(s) \zeta(s)\bigr)
    = 0.
\]
Hence $e^{-L(s)} \zeta(s)$ is constant on $\mathcal{D}$, and
evaluating at $s = c_*$ gives
$e^{-L(c_*)} \zeta(c_*) = e^{-\log\zeta(c_*)} \zeta(c_*) = 1$.
Therefore $e^{L(s)} = \zeta(s)$ throughout $\mathcal{D}$, so $L$
serves as a continuous branch of $\log \zeta$ on $\mathcal{D}$.

We define
\begin{equation}\label{eq:zeta-alpha-def}
    \zeta(s)^\alpha := e^{\alpha L(s)}
    \qquad (s \in \mathcal{D},\ \alpha \in \C).
\end{equation}
This is single-valued and holomorphic in $s$ on $\mathcal{D}$,
and entire in $\alpha$. At $\alpha = 2$, the definition reduces
to $\zeta(s)^2$ in the usual sense.

\subsubsection*{Vertical growth of \texorpdfstring{$\zeta(s)$}{zeta(s)}}

We use the following uniform bound on $|\zeta(s)|$, obtained by
combining the convexity bound with the trivial estimate
$\zeta(\sigma + it) \ll \log |t|$ for $\sigma \geq 1$ (see, e.g.,
\cite[Theorem II.3.8]{} for the convexity bound).

\begin{theorem}[Uniform bound on $\zeta(s)$ in a vertical strip]
\label{thm:zeta-uniform}
Fix $\delta \in (0, 1/2)$. For every $\varepsilon > 0$,
\begin{equation}\label{eq:zeta-uniform}
    |\zeta(\sigma + it)| \ll_{\delta, \varepsilon}
    |t|^{\delta/2 + \varepsilon}
    \quad \text{for } 1 - \delta \leq \sigma \leq 1 + \delta
    \text{ and } |t| \geq 2,
\end{equation}
where the implied constant is independent of $t$.
\end{theorem}

\begin{proof}[Sketch]
For $1 - \delta \leq \sigma \leq 1$, the convexity bound gives
$|\zeta(\sigma + it)| \ll_\varepsilon |t|^{(1-\sigma)/2 + \varepsilon}$.
Taking the supremum of $(1 - \sigma)/2$ over $\sigma \in [1-\delta, 1]$
yields $\delta/2$, so
$|\zeta(\sigma + it)| \ll_\varepsilon |t|^{\delta/2 + \varepsilon}$
on this subrange.

For $1 \leq \sigma \leq 1 + \delta$, the standard estimate
$|\zeta(\sigma + it)| \ll \log |t|$ holds, and
$\log |t| \ll_\varepsilon |t|^\varepsilon \leq |t|^{\delta/2 + \varepsilon}$.
Combining the two ranges gives~\eqref{eq:zeta-uniform}.
\end{proof}

In what follows, we apply Theorem~\ref{thm:zeta-uniform} with the
smaller choice $\delta = \delta_*$ when estimating the growth of
$F(s,y)$ on the nonlocal contour pieces. The larger parameter
$\delta'$ remains in force only for the uniform control of $H(s,y)$.

\subsubsection*{A bound on \texorpdfstring{$|F(s, y)|$}{|F(s,y)|}}

The branch~\eqref{eq:zeta-alpha-def} yields a clean bound on
$|F(s, y)|$ on $\mathcal{D}$. Writing $\alpha = 2 + 2y$ with
$y = \Real(y) + i \Im(y)$,
\[
    |\zeta(s)^\alpha|
    = |e^{\alpha L(s)}|
    = e^{\Real(\alpha L(s))}
    = e^{\Real(\alpha) \Real(L(s)) \,-\, \Im(\alpha) \Im(L(s))}.
\]
Using $e^{L(s)} = \zeta(s)$, we have
$\Real(L(s)) = \log|\zeta(s)|$ and
$\Im(L(s)) = \arg \zeta(s)$ (with the latter referring to the
specific branch determined by $L$). Hence
\[
    |\zeta(s)^\alpha|
    = |\zeta(s)|^{\Real(\alpha)} \cdot e^{-\Im(\alpha) \arg \zeta(s)}.
\]
Since $\Re(\alpha) = 2 + 2\Re(y)$, $\Im(\alpha) = 2\Im(y)$, and
$|\Im(y)| \leq |y| \leq \eta$, we obtain
\begin{equation}\label{eq:F-safe-bound}
    |F(s, y)|
    \leq |H(s, y)| \cdot |\zeta(s)|^{2 + 2\Real(y)}
    \cdot \exp\bigl(2 \eta \, |\arg \zeta(s)|\bigr)
    \qquad (s \in \mathcal{D},\ |y| \leq \eta).
\end{equation}

The factor $|H(s, y)|$ is uniformly bounded on
$\mathcal{R}_{\delta', \eta}$ by
Proposition~\ref{prop:H-analytic}. The factor
$|\zeta(s)|^{2 + 2\Real(y)}$ is controlled by
Theorem~\ref{thm:zeta-uniform}, and the factor
$\exp(2\eta |\arg \zeta(s)|)$ by the following standard
estimate (see, e.g., \cite[\S9.4]{} or
\cite[Theorem~II.3.10]{}).

\begin{theorem}[Bound on $\arg \zeta$]
\label{thm:arg-zeta}
There exists an absolute constant $C_{\arg} > 0$ such that, for
every $s=\sigma+it\in\mathcal{D}$ with $|t| \geq 2$,
\begin{equation}\label{eq:arg-zeta-bound}
    |\arg \zeta(s)| \leq C_{\arg} \log|t|,
\end{equation}
where the branch of $\arg \zeta$ is the one determined by
the logarithm $L$ on $\mathcal{D}$.
\end{theorem}

We now combine these bounds. On the range $|t| \geq 2$,
$1 - \delta_* \leq \sigma \leq 1 + \delta_*$, and $s \in \mathcal{D}$,
applying Theorem~\ref{thm:zeta-uniform} with $\delta = \delta_*$ and
Theorem~\ref{thm:arg-zeta} to~\eqref{eq:F-safe-bound} gives
\begin{equation}\label{eq:F-bound-large-t}
    |F(\sigma + it, y)|
    \ll_{\eta, \varepsilon}
    |t|^{(2 + 2\eta)(\delta_*/2 + \varepsilon)}
    \cdot |t|^{2 \eta C_{\arg}}
    = |t|^{B},
\end{equation}
where
\begin{equation}\label{eq:B-exponent}
    B := B(\eta, \varepsilon)
       := (2 + 2\eta)\!\left(\frac{\delta_*}{2} + \varepsilon\right)
        + 2 \eta C_{\arg}.
\end{equation}
The implied constant in~\eqref{eq:F-bound-large-t} depends on
$\eta$ and $\varepsilon$ but not on $t$ or $y$. Since
$\delta_*<c_0/2$, we may choose $\eta$ and $\varepsilon$ sufficiently
small so that
\begin{equation}\label{eq:B-small}
    B < \frac{c_0}{2}.
\end{equation}
In particular, $B<1$.

On the complementary range $|t| \leq 2$ and
$1 - \delta_* \leq \sigma \leq 1 + \delta_*$, the function $F(s,y)$
is not uniformly bounded on the whole punctured strip, because it
has a singularity at $s=1$. This is harmless for the nonlocal
contour estimates: on the left vertical segment the distance to
$s=1$ is $\lambda_T$, while the local Hankel segment is treated
separately as $I_H$. Thus the following bound is used only on the
nonlocal contour pieces lying in $\mathcal{D}$ and, for all
sufficiently large $x$, inside the narrower strip
$1-\delta_*\leq \sigma \leq 1+\delta_*$:
\begin{equation}\label{eq:F-bound-uniform}
    |F(\sigma + it, y)|
    \ll_{\eta, \varepsilon}
    (|t| + 2)^B
    \qquad (|t|\geq 2,\ |y| \leq \eta).
\end{equation}
For the bounded-height portions of the nonlocal contour, the possible
loss is at most a power of $\log T$ and may be absorbed into the
choice of $\varepsilon$ after $T=T(x)$ is fixed.

\subsection{Contour decomposition}\label{sec:contour-decomposition}

We now construct the deformed contour and use it, together with
the analytic preliminaries of Section~\ref{sec:analytic-prelim},
to express $A(x, y)$ as a sum of contour integrals.

\subsubsection*{The deformed contour}

Fix a truncation height $T \geq 1$, and let $\lambda_T$ be as
in~\eqref{eq:lambda-T}. Set
\begin{equation}\label{eq:sigma0-T}
    \sigma_0 := \sigma_0(T) := 1 - \lambda_T
    = 1 - \frac{c_0}{2 \log(T + 2)}.
\end{equation}
Choose a parameter $\rho > 0$ with $\rho < \lambda_T$ (the precise
choice of $\rho$ as a function of $x$ will be made in
Section~\ref{sec:hankel}). The notation $i0^\pm$ below is a shorthand
for boundary values: more precisely, the two Hankel arms are first
taken at heights $\pm\varepsilon$ with $\varepsilon>0$, and the limit
$\varepsilon\to0^+$ is taken after the integrals are formed. Thus the
actual contours lie in the slit domain $\mathcal{D}$ before taking
boundary values.

We deform the original Perron contour $[c - iT, c + iT]$ leftward,
around the singularity at $s = 1$, into a contour
$\mathcal{C} = \mathcal{C}(T, \rho)$ traversed from $c - iT$ to
$c + iT$ and consisting of the following five pieces:

\begin{itemize}
    \item[(1)] $\mathcal{C}_{\mathrm{bot}}$: from $c - iT$ to
    $\sigma_0 - iT$, leftward;

    \item[(2)] $\mathcal{C}_{\mathrm{left}}^-$: from
    $\sigma_0 - iT$ to $\sigma_0 + i 0^-$, upward;

    \item[(3)] $\mathcal{H}_{\mathrm{loc}}$, a local Hankel contour
    around $s = 1$, consisting of:
    \begin{itemize}
        \item[(3a)] from $\sigma_0 + i 0^-$ to $1 - \rho + i 0^-$,
        rightward along the lower edge of $\mathcal{S}$;

        \item[(3b)] the arc $\{|s - 1| = \rho\}$ from
        $1 - \rho + i 0^-$ to $1 - \rho + i 0^+$, traversed
        counterclockwise (i.e., passing through $s = 1 + \rho$);

        \item[(3c)] from $1 - \rho + i 0^+$ to $\sigma_0 + i 0^+$,
        leftward along the upper edge of $\mathcal{S}$;
    \end{itemize}

    \item[(4)] $\mathcal{C}_{\mathrm{left}}^+$: from
    $\sigma_0 + i 0^+$ to $\sigma_0 + iT$, upward;

    \item[(5)] $\mathcal{C}_{\mathrm{top}}$: from $\sigma_0 + iT$
    to $c + iT$, rightward.
\end{itemize}

The Hankel arms (3a) and (3c) follow the slit \(\mathcal S\) as
upper and lower boundary values. In the local representation
\[
F(s,y)=G(s,y)(s-1)^{-\alpha},
\]
the discontinuity across \(\mathcal S\) is entirely contained in
\((s-1)^{-\alpha}\), and this discontinuity produces the local Hankel
contribution \(I_H\) evaluated in Section~\ref{sec:hankel}.
The five pieces are joined consecutively, so $\mathcal{C}$ begins
at $c - iT$ and ends at $c + iT$, matching the orientation of the
original Perron contour.

\subsubsection*{Verification that the contour avoids \texorpdfstring{$\zeta$}{zeta}-zeros}

We claim that $\mathcal{C} \setminus \{1\} \subset \mathcal{Z}$,
where $\mathcal{Z}$ is the zero-free region of
Theorem~\ref{thm:zero-free}. For the horizontal pieces
$\mathcal{C}_{\mathrm{bot}}$ and $\mathcal{C}_{\mathrm{top}}$,
\begin{equation}\label{eq:zfree-horizontal}
    \Real(s) \geq \sigma_0(T)
    = 1 - \frac{c_0}{2\log(T+2)}
    > 1 - \frac{c_0}{\log(|\Im s| + 2)}.
\end{equation}
For the vertical pieces $\mathcal{C}_{\mathrm{left}}^\pm$,
\begin{equation}\label{eq:zfree-vertical}
    \Real(s) = \sigma_0(T)
    = 1 - \frac{c_0}{2\log(T+2)}
    > 1 - \frac{c_0}{\log(|\Im s| + 2)}.
\end{equation}
For the local Hankel contour $\mathcal{H}_{\mathrm{loc}}$, the two
horizontal arms (3a) and (3c) satisfy $\Real(s) \geq \sigma_0(T)$
and $|\Im s| \leq \rho$, hence are covered by the same estimate
as the vertical pieces. For the small circle (3b),
\begin{equation}\label{eq:zfree-hankel}
    \Real(s) \geq 1 - \rho
    > 1 - \lambda_T
    = 1 - \frac{c_0}{2\log(T+2)}
    > 1 - \frac{c_0}{\log(|\Im s| + 2)}.
\end{equation}

Therefore $\mathcal{C} \setminus \{1\} \subset \mathcal{Z}$.
Combined with the boundary-value convention for the Hankel arms, the
approximating contours with heights $\pm\varepsilon$ are contained in
$\mathcal{D}$ away from $s=1$. The integrand $F(s, y)x^s/s$ is
therefore well-defined and holomorphic on a neighborhood of each
approximating contour, and the contour $\mathcal{C}$ is understood as
the corresponding boundary-value limit.

\subsubsection*{Cauchy's theorem}

The Perron contour and $\mathcal{C}$ share the same endpoints
$c \pm iT$. Together, the Perron contour traversed from $c - iT$
to $c + iT$ and $\mathcal{C}$ traversed in the reverse direction
form a closed loop $\Gamma$, oriented counterclockwise. The
keyhole formed by $\mathcal{H}_{\mathrm{loc}}$ and the slit
$\mathcal{S}$ excludes the point $s = 1$ from the interior of
$\Gamma$, so the interior of $\Gamma$ is contained in the
rectangle
\begin{equation}\label{eq:Gamma-interior-bound}
    \mathcal{R}_\Gamma := \{ s = \sigma + it
    : \sigma_0(T) \leq \sigma \leq c, \ |t| \leq T \}
    \setminus \{1\},
\end{equation}
with the slit further excluded. For every $s \in \mathcal{R}_\Gamma$,
\begin{equation}\label{eq:zfree-interior}
    \Real(s) \geq \sigma_0(T)
    = 1 - \frac{c_0}{2\log(T+2)}
    > 1 - \frac{c_0}{\log(|\Im s| + 2)},
\end{equation}
so $\mathcal{R}_\Gamma \subset \mathcal{Z}$. Combined with the
on-contour verifications~\eqref{eq:zfree-horizontal}--\eqref{eq:zfree-hankel},
this shows that $\Gamma$ together with its interior is contained
in $\mathcal{Z}$. In particular, $F(s, y) x^s / s$ is holomorphic
on a neighborhood of $\Gamma$ and its interior, with the slit
$\mathcal{S}$ removed. By Cauchy's theorem,
\begin{equation}\label{eq:cauchy-zero}
    \oint_\Gamma F(s, y) \frac{x^s}{s}\, ds = 0,
\end{equation}
which yields
\begin{equation}\label{eq:contour-equality}
    \int_{c - iT}^{c + iT} F(s, y) \frac{x^s}{s}\, ds
    = \int_{\mathcal{C}} F(s, y) \frac{x^s}{s}\, ds.
\end{equation}

\subsubsection*{Decomposition of \texorpdfstring{$A(x, y)$}{A(x,y)}}

Combining~\eqref{eq:contour-equality} with the truncated Perron
formula (Proposition~\ref{prop:perron-applied}), we obtain
\begin{equation}\label{eq:A-decomposition}
    A(x, y) = I_H + I_{\mathrm{top}} + I_{\mathrm{bot}}
            + I_{\mathrm{left}} + E_{\mathrm{Perron}}(x, T, y),
\end{equation}
where, for $* \in \{\mathrm{top},\ \mathrm{bot},\ \mathrm{left}\}$,
\begin{equation}\label{eq:I-pieces}
    I_* := \frac{1}{2\pi i} \int_{\mathcal{C}_*}
    F(s, y) \frac{x^s}{s}\, ds,
\end{equation}
with $\mathcal{C}_{\mathrm{left}} := \mathcal{C}_{\mathrm{left}}^-
\cup \mathcal{C}_{\mathrm{left}}^+$, and
\begin{equation}\label{eq:I-H}
    I_H := \frac{1}{2\pi i} \int_{\mathcal{H}_{\mathrm{loc}}}
    F(s, y) \frac{x^s}{s}\, ds.
\end{equation}

The local Hankel contribution $I_H$ contains the main asymptotic
behavior and will be evaluated in Section~\ref{sec:hankel}. The
remaining contributions $I_{\mathrm{top}}, I_{\mathrm{bot}},
I_{\mathrm{left}}$ are estimated in
Section~\ref{sec:contour-estimates} and shown to be of lower
order than the main term, after a suitable choice of $T = T(x)$.

\subsection{Estimates of the contour pieces}\label{sec:contour-estimates}

In this section, we estimate the three nonlocal contour
contributions $I_{\mathrm{top}}, I_{\mathrm{bot}},
I_{\mathrm{left}}$ appearing in the
decomposition~\eqref{eq:A-decomposition}, and combine them with
the Perron truncation error $E_{\mathrm{Perron}}$ to obtain a
single error bound. The local Hankel contribution $I_H$ is
deferred to Section~\ref{sec:hankel}.

The main analytic tool throughout this section is the
bound~\eqref{eq:F-bound-uniform},
\begin{equation*}
    |F(\sigma + it, y)|
    \ll_{\eta, \varepsilon}
    (|t| + 2)^B
    \qquad (1 - \delta_* \leq \sigma \leq 1 + \delta_*,\
    s \in \mathcal{D},\ |t|\geq 2,\ |y| \leq \eta),
\end{equation*}
where $B = B(\eta, \varepsilon)$ is defined
in~\eqref{eq:B-exponent}. We assume throughout that $\eta$ and
$\varepsilon$ are chosen sufficiently small that $B < c_0/2$.

\subsubsection*{Horizontal pieces}

\begin{proposition}\label{prop:horizontal-bound}
For every $\eta \in (0, 1/2)$ and $\varepsilon > 0$ with $B < 1$,
for all sufficiently large $x$ and every sufficiently large $T$ for
which $\lambda_T \leq \delta_*$,
\begin{equation}\label{eq:horizontal-bound}
    |I_{\mathrm{top}}| + |I_{\mathrm{bot}}|
    \ll_{\eta, \varepsilon}
    x^c \, T^{B - 1},
\end{equation}
uniformly in $|y| \leq \eta$.
\end{proposition}

\begin{proof}
On $\mathcal{C}_{\mathrm{top}}$, parameterize $s = \sigma + iT$
with $\sigma_0(T) \leq \sigma \leq c$. Then
$|x^s| = x^\sigma \leq x^c$, $|s| \geq T$, and
$|F(\sigma + iT, y)| \ll_{\eta, \varepsilon} T^B$
by~\eqref{eq:F-bound-uniform} (using $T \geq 2$). The length
$c - \sigma_0(T) \leq c$ is bounded. Therefore
\[
    |I_{\mathrm{top}}|
    \ll_{\eta, \varepsilon}
    \frac{x^c}{T} \cdot T^B
    = x^c \, T^{B - 1}.
\]
The same estimate applies to $|I_{\mathrm{bot}}|$.
\end{proof}

\subsubsection*{Vertical pieces}

\begin{proposition}\label{prop:vertical-bound}
For every $\eta \in (0, 1/2)$ and $\varepsilon > 0$ with $B < c_0/2$,
for all sufficiently large $x$ and every sufficiently large $T$ for
which $\lambda_T \leq \delta_*$,
\begin{equation}\label{eq:vertical-bound}
    |I_{\mathrm{left}}|
    \ll_{\eta, \varepsilon}
    x^{1 - \lambda_T} \, T^B,
\end{equation}
uniformly in $|y| \leq \eta$, where $\lambda_T$ is defined
in~\eqref{eq:lambda-T}.
\end{proposition}

\begin{proof}
The contour $\mathcal{C}_{\mathrm{left}} =
\mathcal{C}_{\mathrm{left}}^- \cup \mathcal{C}_{\mathrm{left}}^+$
parameterizes as $s = \sigma_0(T) + it$ with $|t| \leq T$. We
have $|x^s| = x^{\sigma_0(T)} = x^{1 - \lambda_T}$ and
$|s| \geq \max(\sigma_0(T), |t|) \gg \max(1, |t|)$ for $T$
sufficiently large. Hence
\[
    |I_{\mathrm{left}}|
    \ll_{\eta, \varepsilon}
    x^{1 - \lambda_T}
    \left(
    \int_{2 \leq |t| \leq T} \frac{|t|^B}{|t|}\,dt
    + \int_{|t|\leq 2} \frac{|F(\sigma_0(T)+it,y)|}{1}\,dt
    \right).
\]
For the first integral,
\[
    \int_{2 \leq |t| \leq T} |t|^{B-1}\,dt \ll T^B.
\]
On the bounded-height part $|t|\leq2$, the distance from $s=1$ is
at least $\lambda_T$, and the factorization of $F$ gives
$|F(\sigma_0(T)+it,y)|\ll_{\eta,\varepsilon}(\log T)^{O_\eta(1)}$.
Since every fixed power of $\log T$ is $O_\varepsilon(T^\varepsilon)$,
this contribution is also absorbed into $T^B$ after slightly reducing
the admissible size of $\varepsilon$ in the definition of $B$. This
gives~\eqref{eq:vertical-bound}.
\end{proof}

\subsubsection*{Choice of $T$ and the combined error}

We choose
\begin{equation}\label{eq:T-choice}
    T = T(x) := \exp\!\bigl(\sqrt{\log x}\bigr),
\end{equation}
the standard Selberg--Delange truncation height adapted to the
de la Vall\'ee Poussin zero-free region. With this choice,
$\log T = \sqrt{\log x}$, and
$\lambda_T = c_0 / (2 \sqrt{\log x} + O(1))$, so
\[
    \lambda_T \leq \delta_*,
    \qquad
    c = 1 + \frac{1}{\log x} \leq 1+\delta_*
\]
for all sufficiently large $x$. Hence the nonlocal contour pieces
lie in the narrower strip used in the growth estimates. Moreover,
\begin{equation}\label{eq:x-decay}
    x^{- \lambda_T}
    = \exp(- \lambda_T \log x)
    = \exp\!\left(- \tfrac{c_0}{2} \sqrt{\log x} \cdot (1 + o(1))\right).
\end{equation}

\begin{proposition}\label{prop:nonlocal-error}
With $T = T(x) = \exp(\sqrt{\log x})$, there exists a constant
$c_1 > 0$ (depending on $\eta, \varepsilon$ but not on $x$) such
that
\begin{equation}\label{eq:total-nonlocal-error}
    E_{\mathrm{Perron}}(x, T, y)
    + |I_{\mathrm{top}}| + |I_{\mathrm{bot}}|
    + |I_{\mathrm{left}}|
    \ll_{\eta, \varepsilon}
    x \exp\!\bigl(- c_1 \sqrt{\log x}\bigr),
\end{equation}
uniformly in $|y| \leq \eta$, for all $x$ sufficiently large.
\end{proposition}

\begin{proof}
We bound each term separately and take the worst.

\emph{Perron error.}
By~\eqref{eq:Axy-perron-error} with $T = \exp(\sqrt{\log x})$,
\[
    E_{\mathrm{Perron}}
    \ll_{\eta, \delta'}
    \frac{x (\log x)^{2 + 2\eta}}{T}
    = x (\log x)^{2 + 2\eta} \exp(- \sqrt{\log x}),
\]
which is $\ll x \exp(- (1 - o(1)) \sqrt{\log x})$.

\emph{Horizontal pieces.}
Take the standard choice $c := 1 + 1/\log x$, so that
$x^c = e \cdot x$. By
Proposition~\ref{prop:horizontal-bound},
\[
    |I_{\mathrm{top}}| + |I_{\mathrm{bot}}|
    \ll_{\eta, \varepsilon}
    x \cdot T^{B - 1}
    = x \cdot \exp\!\bigl((B - 1) \sqrt{\log x}\bigr).
\]
Since $B < 1$, this is $\ll x \exp(- (1 - B) \sqrt{\log x})$.

\emph{Vertical piece.}
By Proposition~\ref{prop:vertical-bound}
and~\eqref{eq:x-decay},
\[
    |I_{\mathrm{left}}|
    \ll_{\eta, \varepsilon}
    x \cdot x^{- \lambda_T} \cdot T^B
    = x \cdot \exp\!\bigl(- \tfrac{c_0}{2} \sqrt{\log x}
                          + B \sqrt{\log x}\bigr) (1 + o(1)).
\]
For $B < c_0 / 2$, this is
$\ll x \exp(- (c_0/2 - B) \sqrt{\log x})$.

The constraint $B < c_0 / 2$ is achievable by taking $\eta$ and
$\varepsilon$ sufficiently small, since
$B = (2 + 2\eta)(\delta_*/2 + \varepsilon) + 2\eta C_{\arg}$ tends to
$\delta_*$ as $\eta, \varepsilon \to 0^+$, and $\delta_*$ was chosen
strictly smaller than $c_0/2$ in~\eqref{eq:delta-star}.

Choose $c_1>0$ smaller than each of $1$, $1-B$, and $c_0/2-B$.
For all sufficiently large $x$, the preceding three estimates are then
bounded by the right-hand side of~\eqref{eq:total-nonlocal-error}.
\end{proof}

The right-hand side of~\eqref{eq:total-nonlocal-error} is
$o(x \log x)$, so it will be dominated by the local Hankel
contribution $I_H$ to be analyzed in
Section~\ref{sec:hankel}, which is of order $x \log x$.

\section{Local Hankel analysis}\label{sec:hankel}
% 異뷀썑 ?묒꽦

\section{Extraction of $T_j(x)$ and proof of the main theorem}\label{sec:extraction}
% 異뷀썑 ?묒꽦

\section{Toward $D(x)$: outlook}\label{sec:outlook}
% 異뷀썑 ?묒꽦


% ===== References =====
\begin{thebibliography}{99}

% ?섏쨷???ш린???덊띁?곗뒪 異붽?

\end{thebibliography}

\end{document}
